\chapter{Fronteras y aplicaciones recientes}\label{cap:fronteras}

\begin{panorama}
Cerramos con un recorrido por líneas de investigación activas en las que el cálculo de
Malliavin desempeña un papel central: aprendizaje automático (modelos de difusión
generativos y redes neuronales anchas), métodos numéricos (error débil, multinivel), ecuaciones
estocásticas hacia atrás, ecuaciones de McKean--Vlasov, estadística de difusiones, y
una lista de problemas abiertos.
\end{panorama}

\section{Modelos de difusión generativos}
Los modelos generativos basados en difusiones (\emph{score-based diffusion models}) generan muestras
de una ley de datos $p_0$ invirtiendo en el tiempo una difusión que la convierte en ruido gaussiano.
Si $\dd X_t=b(t,X_t)\,\dd t+\sigma(t)\,\dd W_t$ con $X_0\sim p_0$, el proceso invertido $\bar X_s=X_{T-s}$ satisface
\[
  \dd\bar X_s=\big(-b(T-s,\bar X_s)+\sigma\sigma^\top(T-s)\nabla\log p_{T-s}(\bar X_s)\big)\,\dd s+\sigma(T-s)\,\dd\bar W_s,
\]
y la clave es aprender la \emph{función score} $\nabla\log p_t$. El cálculo de Malliavin da una representación
de la score como esperanza condicional.

\begin{proposicion}[Score por integración por partes]
Sea $X_t\in(\D^\infty)^d$ no degenerado, con
densidad $p_t>0$. Entonces
\[
  \nabla\log p_t(x)=-\E\big[H(X_t)\,\big|\,X_t=x\big],\qquad H_i(X_t)=\sum_j\delta\big((\gamma_{X_t}^{-1})_{ij}DX^j_t\big).
\]
\end{proposicion}

\begin{proof}
Por la Proposición~\ref{prop:ibp-iterada} con $G=1$, $\E[\partial_i\varphi(X_t)]=\E[\varphi(X_t)H_i(X_t)]$. El lado
izquierdo es $\int\partial_i\varphi\,p_t=-\int\varphi\,\partial_ip_t$, y el derecho es $\int\varphi(x)\E[H_i\mid X_t=x]p_t(x)\,\dd x$. Igualando para
toda $\varphi$, $\partial_ip_t=-p_t\E[H_i\mid X_t=x]$.
\end{proof}

Cuando $X_0$ es determinista, esta es la fórmula de Bismut; cuando $X_0\sim p_0$, se aplica
condicionalmente. La consecuencia práctica es que la score es la mejor aproximación en $L^2$ de
$-H(X_t)$ por funciones de $X_t$, de modo que puede aprenderse por regresión (es el análogo de la
\emph{denoising score matching}, que corresponde a elegir pesos de Malliavin particulares). Trabajos de
2025 desarrollan esta conexión para difusiones con coeficientes no lineales y ruido degenerado,
donde las fórmulas de tipo gaussiano no se aplican.

\section{Redes neuronales anchas}
Una red neuronal totalmente conectada con pesos gaussianos aleatorios converge, cuando la
anchura $n\to\infty$, a un proceso gaussiano (Neal; Lee et al.; Matthews et al.). Cada salida de la
red es un funcional no lineal de los pesos gaussianos. Basteri y Trevisan (2022) cuantificaron esta
convergencia con técnicas de transporte óptimo, y Favaro, Hanin, Marinucci, Nourdin y Peccati
(2023) obtuvieron tasas del orden de $n^{-1/2}$ en distancias de Wasserstein y total variación con las
desigualdades de Poincaré de segundo orden del Capítulo~\ref{cap:malliavin-stein}. La derivada de Malliavin respecto
a los pesos de la última capa es, esencialmente, el gradiente que usa la retropropagación.

\section{Métodos numéricos}

\paragraph{Error débil.} Para el esquema de Euler $X^n$ de una EDE, el error débil $\E\Phi(X_T)-\E\Phi(X^n_T)$ es
$O(1/n)$ si $\Phi$ es regular (Talay--Tubaro). Bally y Talay probaron, con integración por partes de
Malliavin y la condición de Hörmander, que la tasa $O(1/n)$ se mantiene para $\Phi$ \emph{solo medible y
acotada}, y que la densidad de $X^n_T$ converge a la de $X_T$ a la misma tasa.

\paragraph{Multinivel.} En el método Monte Carlo multinivel de Giles, los estimadores de griegas con
pesos de Malliavin tienen varianza de nivel $O(1/n)$ incluso para pagos discontinuos, manteniendo
la complejidad óptima $O(\varepsilon^{-2})$ en casos donde la diferenciación trayectorial falla.

\paragraph{Cubatura en el espacio de Wiener.} Kusuoka y Lyons--Victoir construyen medidas discretas
en el espacio de trayectorias que reproducen los momentos de las integrales iteradas del browniano
hasta cierto orden. El análisis del error usa la regularización del semigrupo bajo la condición de
Hörmander, probada con cálculo de Malliavin (condición UFG de Kusuoka).

\section{Ecuaciones estocásticas hacia atrás}
Una EDEA es un par adaptado $(Y,Z)$ con $Y_t=\xi+\int_t^Tf(s,Y_s,Z_s)\,\dd s-\int_t^TZ_s\,\dd W_s$. El resultado
clásico de El Karoui, Peng y Quenez es que, si $\xi$ y $f$ son Malliavin-diferenciables,
\[
  Z_t=D_tY_t\qquad\text{(versión adecuada)},
\]
y $D_rY_t$ satisface a su vez una EDEA lineal. Esto da la regularidad de $Z$, crucial para el análisis
de esquemas de discretización (Zhang; Bouchard--Touzi) y para los métodos de aprendizaje
profundo para EDPs parabólicas no lineales en alta dimensión (Han, Jentzen y E), en los que
$Z$ se aproxima con una red neuronal y las fórmulas de tipo Bismut proporcionan estimadores
alternativos.

\section{Ecuaciones de McKean--Vlasov}
En una ecuación de McKean--Vlasov, $\dd X_t=b(X_t,\mathcal L(X_t))\,\dd t+\sigma(X_t,\mathcal L(X_t))\,\dd W_t$, los coeficientes
dependen de la ley de la solución. El cálculo de Malliavin se combina con la \emph{derivada de Lions}
respecto a la medida: Crisan y McMurray obtuvieron fórmulas de integración por partes y la
regularidad de las densidades bajo condiciones de tipo Hörmander, y se han obtenido fórmulas
de Bismut para el semigrupo sobre el espacio de medidas, útiles en juegos de campo medio.

\section{Estadística de difusiones}
Gobet usó el cálculo de Malliavin para probar la propiedad LAN (normalidad asintótica local)
para difusiones elípticas observadas en tiempo discreto, expresando la derivada del logaritmo de
la verosimilitud de transición, que no es explícita, como una esperanza condicional de pesos
de Malliavin. La misma idea da estimadores de máxima verosimilitud aproximados y cotas de
eficiencia.

\section{Problemas abiertos}
\begin{enumerate}
  \item \textbf{Hörmander para $H\le1/4$.} Para el browniano fraccionario con $H\le1/4$ (o procesos
  gaussianos aún más irregulares), la teoría de rough paths requiere levantamientos de orden
  superior que existen pero no son canónicos, y no se sabe probar la existencia de densidad
  regular en general bajo la condición de Hörmander.
  \item \textbf{Densidades $C^\infty$ para EDPEs singulares} en todo el régimen subcrítico: falta un
  análogo infinito-dimensional robusto de la no degeneración de la matriz de Malliavin
  compatible con la renormalización.
  \item \textbf{Hipoelipticidad en dimensión infinita.} Más allá de los resultados de Hairer--Mattingly
  para EDPEs semilineales con ruido degenerado, el caso de ecuaciones cuasilineales o
  singulares está abierto.
  \item \textbf{Tasas óptimas fuera de un caos fijo.} Para funcionales con infinitos caos (como las
  soluciones de EDPEs), las tasas de las desigualdades de Poincaré de segundo orden no
  son siempre óptimas, y caracterizar la tasa óptima es un problema abierto en general.
  \item \textbf{Malliavin--Stein para límites no gaussianos.} Fuera del segundo caos y de las
  leyes Gamma, la teoría para límites en caos superiores o límites mezcla es incompleta:
  por ejemplo, el teorema de convergencia a variables del segundo caos tiene condiciones
  necesarias y suficientes solo en casos particulares.
  \item \textbf{Ruidos no gaussianos.} Extender el programa de spectral gap del Capítulo~\ref{cap:estructuras} a ruidos
  de Poisson o de Lévy, donde la desigualdad de Poincaré vale pero la hipercontractividad
  no, es un objetivo natural.
\end{enumerate}

\section{Lecturas recomendadas}
\begin{itemize}
  \item \emph{Para la teoría general:} Nualart~\cite{Nualart06} (la referencia estándar), Nualart--Nualart~\cite{NualartNualart18} (más
  accesible).
  \item \emph{Para Malliavin--Stein:} Nourdin--Peccati~\cite{NourdinPeccati12} y la reseña~\cite{AzmoodehPeccatiYang21}.
  \item \emph{Para finanzas:} Alòs--García-Lorite~\cite{AlosGarciaLorite21}, Malliavin--Thalmaier~\cite{MalliavinThalmaier06}, Di Nunno--Øksendal--Proske~\cite{DiNunnoOksendalProske09}.
  \item \emph{Para EDPEs:} Sanz-Solé~\cite{SanzSole05}, Khoshnevisan~\cite{Khoshnevisan14}, Walsh~\cite{Walsh86}.
  \item \emph{Para Hörmander:} Hairer~\cite{Hairer11}, una prueba completa en 15 páginas.
  \item \emph{Para rough paths:} Friz--Hairer~\cite{FrizHairer20}.
  \item \emph{Para estructuras de regularidad:} Broux--Otto--Tempelmayr~\cite{BrouxOttoTempelmayr25}.
\end{itemize}
